\documentclass[aps,prl,onecolumn,superscriptaddress,nofootinbib]{revtex4-2}

\usepackage{amsmath,amssymb,bm}
\usepackage{booktabs}
\usepackage{graphicx}
\usepackage[colorlinks=true,citecolor=blue,linkcolor=blue,urlcolor=blue]{hyperref}

\begin{document}
\renewcommand{\thefigure}{S\arabic{figure}}

\title{Supplemental Material for Geometry-Encoded Hidden Orientational Memory}

\author{Vladimir A. Baulin}
\affiliation{Synthetix Institute, Amsterdam, The Netherlands}
\email{info@synthetix.institute}

\maketitle

\section{Capillary origin of the bond-frame interaction}

At fixed particle geometry and interfacial contact, the interface height
\(u(\bm r)\) obeys the linearized Young--Laplace equation at scales below the
capillary length,
\begin{equation}
 \gamma\nabla^2 u=0.
 \label{eq:supp-young-laplace}
\end{equation}
For an isolated ellipsoid in this regime, the far-field solution is a harmonic
multipole expansion.
Averaging over force- and torque-free modes leaves the leading orientational mode
\begin{equation}
u_i(r,\phi)=\Delta u_{\max}\left(\frac{R}{r}\right)^2\cos2(\phi-\theta_i).
\end{equation}
with \((r,\phi)\) measured from the particle centre and \(\theta_i\) the
apolar axis.  This profile is the static equilibrium shape required to satisfy
anisotropic contact-angle constraints for an ellipsoidal contact line.

For two particles at separation \(r_{ij}\), with \(\phi_{ij}=\arg(\bm r_j-\bm r_i)\) and
\(\varphi_i=\theta_i-\phi_{ij}\), \(\varphi_j=\theta_j-\phi_{ij}\), interface additivity in
this far-field limit gives \(u=u_i+u_j\).  The capillary interaction is the cross
term in the
interface area cost,
\begin{equation}
 U_{ij}^{\rm cap}\approx
 -\gamma\int \nabla u_i\cdot\nabla u_j\,dA
 =-3\pi\gamma(\Delta u_{\max})^2
 \left(\frac{R}{r_{ij}}\right)^4
 \cos\left[2(\varphi_i+\varphi_j)\right],
 \label{eq:supp-capillary-relative}
\end{equation}
where \(\varphi_i,\varphi_j\) are the axes relative to the centre-to-centre
direction.  This is the second term in the
model's pair energy and can be written
\begin{equation}
U_{ij}^{\rm cap}\equiv -A_4(r_{ij})\cos2(\theta_i+\theta_j-2\phi_{ij}),
\qquad
A_4(r_{ij})=3\pi\gamma(\Delta u_{\max})^2\left(\frac{R}{r_{ij}}\right)^4.
\end{equation}
For \(A_4>0\), each isolated bond is minimized by
\(\theta_i+\theta_j=2\phi_{ij}+m\pi\), where \(m\) is an integer.
Different bond angles on a disordered loop generally impose mutually
incompatible minima, so the many-particle ground state is a compromise among
bond-frame constraints rather than a uniform director.
Substitution of the two relative angles gives
\begin{equation}
 \cos(2\varphi_i+2\varphi_j)
 =\cos2(\theta_i+\theta_j-2\phi_{ij}),
 \label{eq:supp-capillary-map}
\end{equation}
which is the second pair harmonic in the Letter.  This identity was tested
numerically to machine precision in
\texttt{tests/test\_rotating\_colloids\_capillary\_pair.py}.  Equation
\eqref{eq:supp-capillary-relative} is a far-field expression.  Near contact,
the interface shape and ellipsoidal geometry generate additional angular
harmonics.  The pairwise approximation is nevertheless accurate for a range
of experimentally studied configurations, and empirical pair potentials can
replace the quadrupolar expression when pair-calibration data are available
\cite{luo2019}.

The first term of the model,
\(-J_{ij}\cos2(\theta_i-\theta_j)\), is the leading apolar relative-angle
harmonic.  It represents short-range steric, near-field capillary, or other
shape-dependent alignment.  Keeping the two harmonics separate allows the
competition between relative alignment and bond-frame alignment to be
measured.  No one-body term of the form
\(-h_i\cos2(\theta_i-\alpha_i)\) is present after the temporary writing
field is removed.

\section{Caged positional graph}

The graph construction starts from an \(n\times n\) triangular monolayer with
nearest-neighbour spacing \(a=1\) and periodic boundaries.  Even values of
\(n\) close the staggered rows in a rectangular periodic cell.  Each centre
receives an independent Gaussian displacement with component standard
deviation \(0.16a\).  Edges connect all pairs within \(2.6a\).  Pair
displacements and bond angles are computed with the minimum-image
convention.  Rare pairs with \(r_{ij}<0.55a\) are rejected to prevent a
divergent \(r^{-4}\) weight.

For each graph, \(r_0\) is the median nearest-neighbour distance.  The
capillary and short-range weights are
\begin{align}
 w^+_{ij}&=\min\left[6,\left(\frac{r_0}{r_{ij}}\right)^4\right],
 \label{eq:supp-cap-weight}\\
 w^-_{ij}&=
 \begin{cases}
 \exp[-\max(r_{ij}/r_0-1,0)/0.2],&r_{ij}\leq1.35r_0,\\
 0,&r_{ij}>1.35r_0.
 \end{cases}
 \label{eq:supp-short-weight}
\end{align}
The cap at 6 is a numerical guard against a rare near-overlap; it is inactive
for the three selected \(N=144\) graphs.  For graph seed 17, the graph has
1654 edges, mean degree 22.97, capillary weighted degree 3.24, and
short-range weighted degree 1.56.  Its bare fourth bond-angle harmonic is
\(\left|\langle e^{4i\phi_{ij}}\rangle_E\right|=0.0040\).  The measured
\(G_2=0.451\) therefore arises from correlations between rotor angles and
bond frames, not from a pre-existing angular bias of the graph.

\section{Rotational dynamics}

All energies are in \(k_{\rm B}T\) and times in \(D_r^{-1}\).  The It\^{o}
equation integrated in the simulation is
\begin{equation}
 d\theta_i=
 -\partial_{\theta_i}\left(\frac{U}{k_{\rm B}T}\right)d(D_rt)
 +\sqrt{2\,d(D_rt)}\,dW_i,
 \label{eq:supp-langevin}
\end{equation}
where the \(dW_i\) are independent Wiener increments.  Euler--Maruyama
integration uses \(D_r\Delta t=0.0025\).  Angles are reduced modulo \(\pi\)
because the rotors are apolar.  For the selected-point controls, each graph
has 12 thermal replicas, 5000 burn-in steps, and 10000 sampled steps, with
observables recorded every 80 steps.  The parameter-map pilot has six
replicas per cell.  The long-time protocols use 12 replicas and extend to
\(D_rt=100\).

\section{Order and memory observables}

The observables \(S,C_2,G_2\), and \(q_{\rm EA}\) are defined in the Letter.
Three additional quantities are used.

The overlap of independent replicas at a common final time is
\begin{equation}
 Q_{ab}=\frac{1}{N}\sum_i e^{2i(\theta_i^{(a)}-\theta_i^{(b)})}.
 \label{eq:supp-replica}
\end{equation}
We plot \(|Q_{ab}|\).  Even independent random angles have a finite
Rayleigh floor
\(\langle|Q_{ab}|\rangle\simeq\sqrt{\pi}/(2\sqrt{N})\); this floor is shown
explicitly in Fig.~S1(c).

For split replicas, two copies start from the same equilibrated configuration
and then receive independent thermal noise.  Their real overlap is
\begin{equation}
 Q_{\rm split}(t)=\frac{1}{N}\sum_i
 \cos2[\theta_i^{(a)}(t)-\theta_i^{(b)}(t)].
 \label{eq:supp-split}
\end{equation}
This protocol measures loss of basin identity without relying on a
time-averaged local mean.

The two-time autocorrelation following a random-state quench is
\begin{equation}
 C_{\rm age}(t_w+\Delta t,t_w)=
 \frac{1}{N}\sum_i
 \left\langle\cos2[\theta_i(t_w+\Delta t)-\theta_i(t_w)]\right\rangle.
 \label{eq:supp-aging}
\end{equation}
Waiting-time dependence at a fixed lag is an aging diagnostic.  It does not,
by itself, establish a thermodynamic glass transition.

\section{Finite-size state diagram and spatial correlations}

The state diagram in the Letter contains 441 coupling cells on a
\(21\times21\) grid.  Every cell is averaged over three independently
disordered \(N=400\) positional graphs.  For each graph, 32 thermal replicas
were equilibrated for 20,000 steps and sampled for 60,000 steps at
\(D_r\Delta t=0.0025\), with observables recorded every 100 steps.  The three
graph rows are averaged before classification.  Each resulting cell is
represented by
\begin{equation}
 \bm x=(S,C_2,G_2,q_{\rm EA},C_{\rm window}).
\end{equation}
The five coordinates are standardized across cells and clustered without
including \(J\) or \(g\).  K-means solutions with \(2\leq k\leq6\) give
silhouette scores \(0.384,0.417,0.451,0.444,\) and \(0.407\),
respectively; \(k=4\) is therefore selected.  Repeating the fit with 20
random initializations gives adjusted Rand index 1.0 relative to the
reported assignment.  Graph sensitivity was tested by exact enumeration of
all 27 ordered bootstrap resamples of the three graph rows.  Every resample
selected \(k=4\); the mean adjusted Rand index relative to the full
assignment was 0.972, with range 0.919--1.000.  Leaving out one graph at a
time gave adjusted Rand index at least 0.955.  The taxonomy is therefore
stable to the sampled positional graphs.  Its precise dimensional
decomposition remains observable-dependent: the silhouette difference
between \(k=4\) and \(k=5\) is 0.0075,
\(\operatorname{corr}(q_{\rm EA},C_{\rm window})=0.986\), and removing
\(C_2\) selects \(k=2\).  We consequently use the clusters to locate a
reproducible memory-bearing population, not to assign a thermodynamic phase
count.  Regime names are assigned only after fitting, from
the observable centroids.  The confidence plotted in the Letter is
\begin{equation}
 m=\frac{d_2-d_1}{d_2},
\end{equation}
where \(d_1\) and \(d_2\) are the distances to the nearest and
second-nearest standardized centroids.  Small \(m\) marks cells for which
the finite-size taxonomy is unresolved.  Neither K-means boundaries nor
their names constitute thermodynamic phase boundaries.

\begin{figure}[t]
\centering
\includegraphics[width=0.98\textwidth]{capillary_prl_figures/capillary_regime_state_space.pdf}
\caption{\label{fig:supp-state-space}
Observable-space projections of the four data-resolved dynamical regimes.
(a) Projection in the \((S,G_2)\) plane.  (b) Projection in the
\((C_2,G_2)\) plane.  (c) Projection in the \((q_{\rm EA},S)\) plane.
Axes are dimensionless order parameters measured from trajectory averages.
Convex envelopes show the extent of each population; they are descriptive
finite-size regions rather than phase boundaries.}
\end{figure}

For the supplementary continuous diagnostic we also use
\(\mathcal H=(1-S)\min(C_2^+,G_2^+,q_{\rm EA})\), where
\(X^+=\max(X,0)\).  Its \(\mathcal H=0.30\) sector is shown only to compare
the cluster-based result with the earlier threshold representation.

For the real-space analysis, three independent \(N=256\) positional graphs
were simulated at \(J=4\) and \(g=5\), with eight thermal replicas, 8000
burn-in steps, and 12000 sampled steps.  Orientations were recorded every
100 steps, giving 960 configurations per graph.  For every particle pair,
we measured
\begin{align}
 C_2(r)&=\left\langle\cos2(\theta_i-\theta_j)\right\rangle_{r_{ij}\in r},
 \\
 G_2(r)&=\left\langle
 \cos2(\theta_i+\theta_j-2\phi_{ij})\right\rangle_{r_{ij}\in r}.
\end{align}
The connected curve in Fig.~\ref{fig:supp-internal}(b) subtracts the sampled
\(\langle S^2\rangle=0.0144\) from \(C_2(r)\).  No parametric decay law is
fitted.  We define the operational correlation range as the first bin after
the first-shell maximum whose magnitude lies below three times the
large-\(r\) tail noise.  Both angular channels meet this criterion by
\(r/a=1.73\).  The analysis is generated by
\texttt{scripts/analyze\_rotating\_colloids\_capillary\_internal.py}.

\begin{figure}[t]
\centering
\includegraphics[width=0.98\textwidth]{capillary_prl_figures/fig4_capillary_internal_structure.pdf}
\caption{\label{fig:supp-internal}
Internal diagnostics underlying the observable-space phase classification.
(a) At fixed \(J=4k_{\rm B}T\), capillary coupling raises \(G_2\) and
\(q_{\rm EA}\), suppresses \(S\), and leaves finite \(C_2\).  Each bar pair
gives a dimensionless observable for one coupling point.
(b) Connected relative-angle and bond-frame correlations averaged over
three \(N=256\) positional graphs; bands are graph-to-graph standard
deviations.  The correlations are restricted to the first few coordination
shells.  (c) The 1323 graph-cell observations spanning 441 coupling cells in
the \((S,q_{\rm EA})\) projection; the star marks the dynamical test point.}
\end{figure}

\section{Write-release-read protocol}

A target configuration is first obtained from an independently equilibrated
copy of the autonomous pair model.  Independent replicas are then exposed
temporarily to the conjugate field
\begin{equation}
 U_{\rm write}/k_{\rm B}T=-h\sum_i
 \cos2(\theta_i-\theta_i^{\rm target}),
 \qquad h=1.5.
 \label{eq:supp-write}
\end{equation}
The field is removed completely before the release trajectory.  Readout uses
\begin{equation}
 Q_{\rm target}(t)=\frac{1}{N}\sum_i
 \cos2[\theta_i(t)-\theta_i^{\rm target}].
 \label{eq:supp-target}
\end{equation}
At \(N=144,J=4,g=5\), the overlap at removal is 0.801 at the first sampled
release point and remains 0.457 after \(D_rt=99.8\).  The corresponding
split-replica overlap is 0.529.  The field therefore selects a pre-existing
long-lived basin; it is not a permanent substrate term.  Repeating the same
write protocol with \(g=0\) gives an initial release overlap of
\(0.807\pm0.005\), but it falls below 0.2 by
\(D_r(t-t_{\rm off})=4.3\pm0.3\) after release and reaches
\(0.006\pm0.004\) at \(D_rt=625\).

\section{Matched controls}

\begin{table}[ht]
\centering
\caption{Selected-point observables.  Means and standard deviations are over
five quenched positional graphs.  All rows use \(N=1024\) and \(J=4\).
The regular and shuffled controls retain \(g=5\).}
\label{tab:supp-controls}
\begin{ruledtabular}
\begin{tabular}{lccccc}
control & \(S\) & \(C_2\) & \(G_2\) & \(q_{\rm EA}\) &
window correlation\\
\hline
disordered capillary &
\(0.0665\pm0.0038\)&\(0.627\pm0.009\)&\(0.447\pm0.011\)&
\(0.658\pm0.007\)&\(0.543\pm0.007\)\\
same graph, \(g=0\) &
\(0.131\pm0.010\)&\(0.813\pm0.003\)&\(0.0000\pm0.0002\)&
\(0.0225\pm0.0023\)&\(0.0012\pm0.0076\)\\
regular lattice &
\(0.945\pm0.0001\)&\(0.944\pm0.0000\)&\(0.0054\pm0.0000\)&
\(0.818\pm0.005\)&\(0.711\pm0.048\)\\
shuffled bond frames &
\(0.0555\pm0.0055\)&\(0.596\pm0.013\)&\(0.532\pm0.008\)&
\(0.748\pm0.008\)&\(0.669\pm0.013\)\\
\end{tabular}
\end{ruledtabular}
\end{table}

The regular lattice demonstrates that high persistence alone does not imply
hidden memory: it has a unique macroscopic director and \(S\simeq0.96\).
The \(g=0\) graph demonstrates that disorder plus first-shell alignment does
not account for the observed persistence.  Shuffling \(\phi_{ij}\) among
fixed edges preserves the one-edge angle and weight distributions but
destroys their geometric embedding.  Its stronger persistence is consistent
with increased random-bond frustration.  It is a useful upper-frustration
control, not an experimental capillary graph.

\section{Finite-size behavior}

\begin{table}[ht]
\centering
\caption{Finite-size results at \(J=4,g=5\).  Values are means over five
quenched graphs.}
\label{tab:supp-size}
\begin{ruledtabular}
\begin{tabular}{rccccc}
\(N\)&\(S\)&\(C_2\)&\(G_2\)&\(q_{\rm EA}\)&\(|Q_{ab}|\)\\
\hline
144&0.151&0.628&0.445&0.629&0.180\\
256&0.125&0.640&0.442&0.643&0.136\\
576&0.088&0.635&0.443&0.674&0.095\\
1024&0.066&0.627&0.447&0.658&0.075\\
2304&0.040&0.626&0.446&0.647&0.048\\
\end{tabular}
\end{ruledtabular}
\end{table}

Power-law fits over the five sizes give
\(S\propto N^{-0.479}\) and
\(\langle|Q_{ab}|\rangle\propto N^{-0.466}\).  Graph-level bootstrap
resampling gives 95\% intervals
\([-0.510,-0.451]\) and \([-0.477,-0.454]\), respectively.  In contrast,
the exponent of \(q_{\rm EA}\) is 0.011 with interval
\([-0.005,0.039]\); the intervals for \(C_2\) and \(G_2\) likewise contain
zero.  Global order and independent-replica overlap therefore approach their
random finite-size floors, while local correlations and finite-window
persistence remain of order unity.  The largest calculation contains
\(N=2304\) rotors, with five positional-disorder realizations at every size.

\begin{figure}[t]
\centering
\includegraphics[width=0.98\textwidth]{capillary_prl_figures/figS1_capillary_controls_scaling.pdf}
\caption{\label{fig:supp-controls}
Matched controls and finite-size tests at \(J=4\) and \(g=5\).
(a) The physical disordered capillary graphs combine low global order with
large finite-window memory.  Here the horizontal axis is the coupling scan
index and the vertical axis is each measured observable.  Removing \(g\)
removes the memory, whereas the regular lattice selects a global nematic.
Shuffled bond frames are a strong-frustration control that no longer has a
capillary embedding.
(b) \(S\) decreases from \(N=144\) to 2304 while \(C_2\), \(G_2\), and
\(q_{\rm EA}\) remain stable.  Error bars are standard deviations over five
quenched positional graphs.  (c) Independent replicas select different
states; their overlap approaches the random \(N^{-1/2}\) finite-size floor.}
\end{figure}

\section{Equilibrium spin-glass discriminant}

\begin{figure*}[t]
\centering
\includegraphics[width=0.98\textwidth]{capillary_prl_figures/fig4_capillary_spin_glass_discriminant.pdf}
\caption{\label{fig:supp-spin-glass}
Finite-size discriminant between dynamical memory and equilibrium glass
order.  Each point averages five quenched graphs.  (a) The finite-window
single-trajectory statistic \(q_{\rm EA}\) grows along
\((J,g)=\lambda(4,5)k_{\rm B}T\), where \(\lambda\) is the horizontal axis.
(b) The dimensionless overlap correlation length \(\xi_L/L\) is shown as a
function of size and has no common crossing.
(c) Independently equilibrated replicas obey
\(\langle|Q_{ab}|\rangle\propto N^{-1/2}\); lines are log--log fits.
(d) The overlap susceptibility \(\chi_Q=N\langle|Q_{ab}|^2\rangle\)
saturates.}
\end{figure*}

To distinguish finite-window persistence from equilibrium replica ordering,
we performed a separate 200-point scan along
\begin{equation}
 (J,g)=\lambda(4,5)k_{\rm B}T,
 \qquad
 \lambda\in\{0.30,0.45,0.60,0.75,0.90,1.05,1.20,1.40\}.
 \label{eq:supp-ray}
\end{equation}
The scan used \(N=144,256,576,1024,2304\), five independent quenched graphs
at every point, and 48 independently initialized thermal replicas per graph.
Each replica was equilibrated for 50,000 steps and sampled for 100,000 steps.

For each replica pair define the local complex overlap
\(q_i^{ab}=\exp[2i(\theta_i^a-\theta_i^b)]\) and
\(Q_{ab}=N^{-1}\sum_iq_i^{ab}\).  The zero-wave-vector overlap
susceptibility is
\begin{equation}
 \chi_Q=N\langle|Q_{ab}|^2\rangle.
 \label{eq:supp-chiq}
\end{equation}
Writing \(\chi_Q(\bm k)=N\langle|N^{-1}\sum_iq_i^{ab}
e^{i\bm k\cdot\bm r_i}|^2\rangle\), the directional second-moment lengths are
\begin{equation}
 \xi_\mu=\frac{1}{2\sin(k_{\min,\mu}/2)}
 \left[\frac{\chi_Q(0)}{\chi_Q(k_{\min,\mu})}-1\right]^{1/2},
 \qquad \mu=x,y,
 \label{eq:supp-xi}
\end{equation}
and we report \(\xi_L=(\xi_x\xi_y)^{1/2}\) divided by the geometric-mean box
length.  A finite-temperature glass transition would normally produce a
size-independent crossing of \(\xi_L/L\), together with nonvanishing
inter-replica overlap or a growing \(\chi_Q\).

No such signature is present.  At \(\lambda=0.30,0.75,1.40\), log--log fits
give
\begin{equation}
 \langle|Q_{ab}|\rangle\propto
 N^{-0.490},\quad N^{-0.472},\quad N^{-0.485},
 \label{eq:supp-overlap-exponents}
\end{equation}
respectively.  The corresponding largest-size susceptibilities are 1.74,
7.10, and 7.78; they remain finite as \(N\) grows.  The \(\xi_L/L\) curves
decrease with size and do not share a crossing over the scanned interval.
Thus the large finite-window \(q_{\rm EA}\) measures slow trajectory memory,
not an equilibrium Edwards--Anderson order parameter.  The distinction is
important: the write--release and split-replica protocols demonstrate
retention, while independent replicas demonstrate the absence of a common
frozen state.

\section{Programmable mosaic easy-axis realization}

The autonomous model in the Letter stores its local reference frames in the
centre-to-centre bonds.  We also simulated a programmable construction in
which a surface pattern supplies prescribed easy axes \(\alpha_i\).  In this
section only, the dimensionless angular potential is
\begin{equation}
 \Phi_{\rm gr}=-j\sum_{(ij)}w_{ij}\cos2(\theta_i-\theta_j)
 -h\sum_i\cos2(\theta_i-\alpha_i),
 \label{eq:supp-groove-potential}
\end{equation}
and the dynamics obey
\begin{equation}
 \partial_\tau P=\sum_i\partial_{\theta_i}
 \left(\partial_{\theta_i}P+P\partial_{\theta_i}\Phi_{\rm gr}\right),
 \qquad \tau=D_rt.
 \label{eq:supp-groove-smol}
\end{equation}
Here \(j\) and \(h\) are deterministic angular-relaxation rates normalized by
rotational diffusion.  The first term transmits orientation
between neighbours; the second registers each rotor to the local patterned
axis.  Its direct readout is
\begin{equation}
 G_2^{\alpha}=N^{-1}\sum_i\cos2(\theta_i-\alpha_i).
 \label{eq:supp-groove-registration}
\end{equation}

The archived construction is deliberately segmented.  The static scan uses
a \(32\times32\) array divided into sixteen \(8\times8\) patches; the dynamic
protocol uses a \(16\times16\) array with four patches.  Patch axes advance by
\(\pi/4\), so their laboratory-frame nematic contributions cancel.  Within a
patch the interaction graph is dense, whereas adjacent patches are joined by
two links of weight 0.18.  The results therefore test a patterned
easy-axis device with weak interpatch coupling, not a homogeneous grooved
monolayer.

The static \(29\times29\) scan contains 841 parameter cells and 48 thermal
replicas per cell.  The operational conditions
\(S\leq0.35\), \(C_2\geq0.70\), \(G_2^{\alpha}\geq0.70\), and
\(q_{\rm EA}\geq0.50\) select 824 cells in one connected grid component.  A
stricter cut, \(S\leq0.05\), \(C_2,G_2^{\alpha}\geq0.80\), and
\(q_{\rm EA}\geq0.70\), retains 688 cells.  These counts describe correlated
grid cells, not independent samples or a thermodynamic phase.  They establish
that low global order and strong pattern registration occupy an extended
parameter region of this finite device model.

Static persistence is not an independent memory test here.  Across the
mosaic scan, \(q_{\rm EA}\) and \((G_2^{\alpha})^2\) have correlation
0.999987 and root-mean-square difference 0.00178.  This near identity is
expected when each stationary angular distribution is symmetric around its
prescribed axis.  It shows that the field-on statistic records the imposed
pattern.  Removal and rewriting protocols are needed to measure retention.

\begin{figure}[t]
\centering
\includegraphics[width=0.98\textwidth]{grooved_prl_figures/fig2_static_state_space.png}
\caption{\label{fig:supp-groove-static}
Static response of the programmable mosaic model.  (a) Global nematic order
versus registration to the prescribed local axes.  (b) Neighbour coherence
versus field-on persistence.  Curves are smoothed display densities and open
symbols are construction centroids; the simulated points themselves are not
displaced.  (c) The orientational correlation changes sign near the
eight-particle patch width, confirming that the registered pattern follows
the imposed mosaic rather than one common director.  The uniform,
long-range, and triangular scans have different graph sizes and parameter
ranges and are shown as descriptive controls, not as matched statistical
ensembles.}
\end{figure}

The dynamical protocols use 16 replicas.  At \(h=0.2\), raising \(j\) from 0
to 3 increases \(G_2^{\alpha}\) from 0.0965 to 0.8242, while \(S=0.128\).
Neighbour coupling therefore amplifies a weak local cue by a factor of 8.55
without selecting one macroscopic director.  A stronger field, \(h=0.6\), is
then used to write pattern \(A\).  After complete removal of that field, the
overlap with the written configuration has half-lives
\begin{equation}
 \tau_{1/2}=0.177,\ 3.816,\ 5.247
 \quad\text{for}\quad j=0,\ 2,\ 3,
 \label{eq:supp-groove-release}
\end{equation}
corresponding to 22-fold and 30-fold retention enhancements for the two
interacting cases.  The overlap eventually vanishes, and the high-\(j\)
systems develop a conventional director at long times after the patterned
field is absent.  The stored state is therefore volatile.

For rewriting, pattern \(A\) is replaced at \(\tau=0\) by an incompatible
pattern \(B\), with
\(N^{-1}\sum_i\cos2(\alpha_i^A-\alpha_i^B)=-1\).  The time at which the old
overlap falls below 0.5 increases from 0.156 at \(j=0\) to 0.590 and 0.730 at
\(j=2\) and 3.  At long times the system registers to \(B\).  Coupling thus
controls both retention after removal and resistance to replacement.

\begin{figure}[t]
\centering
\includegraphics[width=0.98\textwidth]{grooved_prl_figures/fig3_hidden_memory_evidence.png}
\caption{\label{fig:supp-groove-dynamics}
Dynamic tests of the programmable mosaic.  (a) Survival functions of the
descriptive static score for the four archived constructions.  (b)
Field-supported angular self-correlation at \(h=0.2\); dotted lines show the
corresponding finite-window values.  (c) Overlap with pattern \(A\) after the
writing field is set to zero.  Open symbols mark half-lives and the dotted
line is the free-rotor result.  (d) Overlap with \(A\) after switching to an
incompatible pattern \(B\); symbols mark rewrite half-times.  Curves in
(b--d) average 16 replicas and bands show one standard deviation.}
\end{figure}

Permanent microgrooves can test the static registration and amplification in
Fig.~\ref{fig:supp-groove-static}.  Field removal and pattern replacement
require a switchable implementation, such as a photoresponsive anisotropic
layer or a projected polarization pattern.  Digital-micromirror
photoalignment has produced arbitrary two-dimensional easy-axis fields
\cite{wu2012}, and polarization-resolved optical measurements can recover
individual nanorod orientations \cite{chang2010}.  For a colloidal test, the
single-particle distribution \(p(\theta-\alpha)\) first determines \(h\), and
a two-particle measurement within one patch determines \(j\).  A surface
matching the simulation must also reduce orientational coupling across patch
boundaries, for example by increasing the interpatch spacing.  This extra
fabrication requirement is why the programmable construction is presented as
a Supplemental realization rather than the primary autonomous claim.

\section{Separate weak-link mosaic control}

An earlier groove-free mosaic construction rotated locally square domains and joined
them with a fixed number of weak crosslinks.  It produced nearly simultaneous
large \(C_2\) and \(G_2\), but only 0.6228\% of its total coupling weight
crossed domain boundaries.  In the strong-locking limit its domain directors
reduce exactly to an Ising Hamiltonian
\begin{equation}
 H_{\rm eff}/J=-\sum_{bc}K_{bc}s_bs_c,\qquad s_b=\pm1.
\end{equation}
Exact enumeration found 16--78 single-domain-flip-stable states across 20
graph jitters, confirming a multibasin mechanism.  The same calculation also
showed that the ground-state unsatisfied weight was only 0.75\% of the total
crosslink weight.  The mosaic model is therefore retained as a mechanism
test, not as evidence for the physical capillary regime.  The Letter instead
uses a caged disordered monolayer with no imposed orientational domains.

\section{Experimental parameter estimate}

The required separation between positional and rotational motion has an
experimental precedent.  Charged silica rods have been observed in a state
with strongly caged translational coordinates and nearly free rotations
\cite{besseling2026}; DNA-based sliding contacts provide a second method for
constructing thermally driven colloidal pivots \cite{melio2026}.  Neither
experiment implements the capillary memory interaction, but both establish
that a persistent positional scaffold can coexist with Brownian angular
motion.

To close the static-torque assumption, we additionally require the interfacial
deformation to remain quasistatic over the rotational time scale.
For a meniscus mode of lateral extent \(\ell\), a conservative viscous estimate is
\[
 \tau_{\rm int} \sim \frac{\eta_{\rm eff}\,\ell}{\gamma},
\qquad
 \tau_{\rm int}D_r \ll 1,
\]
with \(\eta_{\rm eff}\) the effective (sum) interfacial viscosity and \(\gamma\) the
interfacial tension.  With \(\eta_{\rm eff}=1\)--\(5\,\mathrm{mPa\,s}\),
\(\gamma=10\,\mathrm{mN\,m^{-1}}\), and \(\ell\sim 2R=4\)--\(20\,\mu\mathrm m\),
this gives \(\tau_{\rm int}\sim10^{-7}\)--\(10^{-5}\,\mathrm s\), so the condition is
safely met when \(D_r^{-1}\) is at least \(10^{-2}\,\mathrm s\) and is typically
much longer in the reported caged interfacial protocol.
The same check is implemented by
\texttt{scripts/estimate\_capillary\_timescale.py}.

For the interfacial realization, the pair torque should be measured before a
dense layer is assembled.  Two ellipsoids held at separation \(r\) give
\begin{equation}
 \frac{U_{\rm eff}(\theta_i,\theta_j\mid r)}{k_{\rm B}T}
 =-\ln p(\theta_i,\theta_j\mid r)+\mathrm{const.}
 \label{eq:supp-pair-calibration}
\end{equation}
from their stationary angular histogram.  Projecting
Eq.~\eqref{eq:supp-pair-calibration} onto
\(\cos2(\theta_i-\theta_j)\) and
\(\cos2(\theta_i+\theta_j-2\phi_{ij})\) determines the two coefficients in
the model.  Repeating the measurement at several separations tests the
\(r^{-4}\) capillary amplitude independently of collective memory.

Solving Eq.~(1) of the Letter for the meniscus amplitude gives
\begin{equation}
 \Delta u_{\max}=
 \left[
 \frac{A_4}{3\pi\gamma}
 \left(\frac{r}{R}\right)^4
 \right]^{1/2}.
 \label{eq:supp-meniscus}
\end{equation}
At room temperature, take \(\gamma=10\,\mathrm{mN\,m^{-1}}\) and \(r/R=3\).
The corresponding energy--amplitude pairs are
\begin{equation}
 \left(\frac{A_4}{k_{\rm B}T},\frac{\Delta u_{\max}}{\mathrm{nm}}\right)
 =(1,1.88),\ (3,3.26),\ (5,4.20).
\end{equation}
Smooth
micron-scale ellipsoids often generate much larger deformations and
interactions far above \(k_{\rm B}T\) \cite{loudet2005}.  Surface
nanostructure can reduce the attraction by an order of magnitude and allow
electrostatic repulsion to maintain finite particle separations
\cite{trevenen2023}.  However, strongly porous ellipsoids lose the
quadrupolar meniscus.  A suitable experiment should therefore tune roughness,
wetting, charge, salt concentration, and surface coverage jointly, and verify
the pair torque before attempting the many-particle memory protocol.

\section{Curvature writing and wave gating}

A deformation of the host interface couples directly to the
particle-induced quadrupole.  In the small-slope far field, the orienting
energy can be written
\begin{equation}
 U_i^{\rm curv}=-\pi\gamma H_pR_p^2\Delta c(\bm r_i,t)
 \cos2(\theta_i-\alpha_i),
 \label{eq:supp-curvature-write}
\end{equation}
where \(\Delta c\) is the local deviatoric curvature, \(\alpha_i\) is its
principal axis, \(H_p\) is the particle-induced quadrupole amplitude, and
\(R_p\) is its lateral scale.  This coupling has been measured through
curvature-driven orientation and migration of interfacial rods
\cite{cavallaro2011}.  It provides a native write protocol: apply a shaped
curvature field, allow the angular state to relax, flatten the interface, and
measure \(Q_{\rm target}\) only after Eq.~\eqref{eq:supp-curvature-write}
has vanished.

For \(\gamma=10\,\mathrm{mN\,m^{-1}}\), \(H_p=5\,\mathrm{nm}\), and
\(R_p=1\,\mu\mathrm m\), a \(k_{\rm B}T\) orienting scale requires
\(\Delta c\simeq26\,\mathrm{m^{-1}}\).  If the imposed interface profile is
\(h=A\cos kx\), this curvature corresponds to
\(A\simeq6.6\,\mathrm{nm}\) at wavelength \(100\,\mu\mathrm m\), or
\(A\simeq0.66\,\mu\mathrm m\) at \(1\,\mathrm{mm}\).  These values are
order-of-magnitude targets because the numerical prefactor depends on the
wetting geometry and curvature convention.

A spatially uniform oscillation does not choose \(\alpha_i\), but it can
modulate the pair barriers.  With
\(\Delta u(t)=\Delta u_0[1+\epsilon\cos\omega t]\) and
\(g\propto\gamma(\Delta u)^2\),
\begin{equation}
 \frac{g(t)}{g_0}\simeq
 1+2\epsilon\cos\omega t+
 \frac{\epsilon^2}{2}\left(1+\cos2\omega t\right).
 \label{eq:supp-wave-gate}
\end{equation}
Equation~\eqref{eq:supp-wave-gate} predicts an interaction gate: weak forcing
anneals local angular basins, whereas a trough that crosses the onset seen in
Fig.~4 of the Letter accelerates erasure.  Low-frequency capillary waves
have annealed air--water colloidal monolayers \cite{caso2025}, and stronger
ultrasonic forcing produces dynamic capillary interactions together with
hydrodynamic rearrangement \cite{huerre2018}.  A memory experiment must
therefore record both \(\theta_i(t)\) and \(\bm r_i(t)\).  The useful
separation of scales is
\(\tau_{\rm int}\ll\omega^{-1}\sim D_r^{-1}\ll\tau_{\rm cage}\): the
interface follows the drive, orientations respond within a cycle, and the
constraint graph remains fixed.

The many-particle experiment then has a direct rejection sequence.  First,
verify \(\tau_{\rm cage}D_r\gg1\) from position and orientation trajectories.
Second, measure \(S,C_2,G_2\) without a writing field.  Third, split repeated
preparations by applying independent thermal histories, or write a target
with a temporary polarized optical field and read only after removal.  The
claim is rejected if low \(S\) is accompanied by vanishing \(G_2\), if the
field-free overlap is indistinguishable from the low-quadrupole control, or
if retention disappears when the observation window is extended.

\section{Discovery and validation sequence}

The calculations were organized so that the same statistic did not both
define and validate the result.  The exploratory stage was the 441-cell
observable-space map.  Coupling coordinates and physical names were excluded
from clustering.  The selected point \((J,g)=(4,5)k_{\rm B}T\) was then fixed
before the larger-size controls and long trajectories were analysed.  The
confirmatory stages were: five-graph selected-point replication, removal of
the capillary channel, regular and shuffled-frame controls, split-replica and
write--release dynamics, and the independent 200-point finite-size overlap
scan.  The programmable groove calculations use a different Hamiltonian and
are not pooled with these tests; they show how the same readout logic can be
implemented with an externally specified angular code.

\section{Publication-scale run}

The complete GPU validation is encoded in
\texttt{scripts/run\_rotating\_colloids\_capillary\_pair\_prl\_gpu.sh}.
It performs:
\begin{enumerate}
\item a \(21\times21\) state map at \(N=400\), with three quenched graphs
and 32 thermal replicas per graph;
\item selected-point and \(g=0\) controls for
 \(N=144,256,576,1024,2304\) across five quenched graphs;
\item regular and shuffled-frame matched controls at \(N=1024\); and
\item overlap, aging, and write-release-read trajectories to
 \(D_rt=625\) for three \(N=1024\) graphs; and
\item an eight-coupling retention scan at \(N=576\), with five quenched
graphs and 32 thermal replicas per graph, used for Fig.~4 of the Letter.
\end{enumerate}
The simulation point files are append-only JSONL records.  Restarting the
same command resumes completed cells.

The final scan is encoded in
\texttt{scripts/run\_rotating\_colloids\_activated\_memory\_prl\_4gpu.sh}.
Each graph is equilibrated for 50,000 steps and observed for 250,000 steps at
\(D_r\Delta t=0.0025\). Its source record is
\texttt{activated\_memory\_scan.jsonl}. Figure~4(b) reports
\begin{equation}
 \mathcal A_Q(T_{\rm obs})=
 \frac{\int_0^{T_{\rm obs}}\max[Q(t),0]dt}{|Q(0)|},
 \label{eq:supp-integrated-overlap}
\end{equation}
a finite-window positive integrated overlap. This quantity compares retained
overlap over the common observation window; it is not an intrinsic
relaxation lifetime and assumes no exponential decay law.

The independent finite-size glass test is encoded in
\texttt{scripts/run\_rotating\_colloids\_spin\_glass\_prl\_4gpu.sh}; its
analysis is generated by
\texttt{scripts/analyze\_rotating\_colloids\_spin\_glass\_prl.py}.

The programmable easy-axis protocols are encoded in
\path{scripts/rotating_colloids_groove_protocols.py}; the two figures are
generated from the archived JSON and JSONL records by
\path{scripts/plot_rotating_colloids_groove_evidence.py}.

\begin{thebibliography}{10}

\bibitem{luo2019}
A. M. Luo, J. Vermant, P. Ilg, Z. Zhang, and L. M. C. Sagis,
J. Colloid Interface Sci. \textbf{534}, 205 (2019).

\bibitem{wu2012}
H. Wu \textit{et al.},
Opt. Express \textbf{20}, 16684 (2012).

\bibitem{chang2010}
W.-S. Chang, J. W. Ha, L. S. Slaughter, and S. Link,
Proc. Natl. Acad. Sci. USA \textbf{107}, 2781 (2010).

\bibitem{besseling2026}
T. H. Besseling \textit{et al.},
Nat. Commun. \textbf{17}, 2410 (2026).

\bibitem{melio2026}
J. Melio, M. van Hecke, S. E. Henkes, \textit{et al.},
Nature \textbf{651}, 632 (2026).

\bibitem{loudet2005}
J.-C. Loudet, A. M. Alsayed, J. Zhang, and A. G. Yodh,
Phys. Rev. Lett. \textbf{94}, 018301 (2005).

\bibitem{trevenen2023}
S. Trevenen \textit{et al.},
ACS Nano \textbf{17}, 11892 (2023).

\bibitem{cavallaro2011}
M. Cavallaro, Jr., L. Botto, E. P. Lewandowski, M. Wang, and K. J. Stebe,
Proc. Natl. Acad. Sci. USA \textbf{108}, 20923 (2011).

\bibitem{caso2025}
M. J. Caso \textit{et al.},
Langmuir \textbf{41}, 3033 (2025).

\bibitem{huerre2018}
A. Huerre, M. De Corato, and V. Garbin,
Nat. Commun. \textbf{9}, 3620 (2018).

\end{thebibliography}

\end{document}
